% =============================================================================
\section{Legacy vs Graph-Based Paradigm: the CH2/CCH2 Coupled System}
\label{sec:comparison}
% =============================================================================

This section contrasts the \emph{legacy} single-tank-first approach with the
\emph{graph-based} multi-tank framework using the coupled CH2/CCH2 system as a
concrete example.  The two approaches produce identical governing equations; the
difference lies entirely in how those equations are constructed, parameterised,
and extended.

% -----------------------------------------------------------------------------
\subsection{Case Description}
\label{sec:comparison:case}
% -----------------------------------------------------------------------------

The system consists of two onboard hydrogen tanks serving an ATR~72 regional
aircraft via a pressure-compensation topology:

\begin{center}
\begin{tabularx}{\linewidth}{@{} l l l l @{}}
\toprule
& \textbf{Tank 1 --- CH2} & \textbf{Tank 2 --- CCH2} \\
\midrule
Fluid          & Compressed GH2  & Cryo-compressed H2 \\
$P_\text{init}$  & \SI{800}{\bar}  & \SI{400}{\bar} \\
$T_\text{init}$  & \SI{288}{\kelvin} & \SI{53.25}{\kelvin} \\
$P_\text{min}$   & \SI{15}{\bar}   & \SI{15}{\bar} \\
$P_\text{vent}$  & \SI{850}{\bar}  & \SI{450}{\bar} \\
Mission         & --- (reserve)      & ATR~72 discharge \\
\bottomrule
\end{tabularx}
\end{center}

\noindent Tank~1 acts as a high-pressure reserve.  A pressure-compensation
valve transfers CH2 from Tank~1 to Tank~2 whenever Tank~2's pressure falls
below \SI{16}{\bar}, preventing supply interruption.  The orifice diameter
is \SI{5}{\milli\metre} and the maximum allowable coupling flow rate is
\SI{56}{\gram\per\second}.

% -----------------------------------------------------------------------------
\subsection{Legacy Paradigm: Monolithic Per-Scenario ODE}
\label{sec:comparison:legacy}
% -----------------------------------------------------------------------------

In the legacy framework, analysing any multi-tank system required the developer
to define a \emph{monolithic} ODE right-hand-side function tailored to that
specific topology.  For the CH2/CCH2 case the developer would:

\begin{enumerate}[label=\arabic*.]
  \item \textbf{Manually define the state vector} and the index convention:
    \begin{equation}
      \vect{y} = \bigl[\,m_1,\; T_1,\; T_{s,1},\; m_2,\; T_2,\; T_{s,2}\,\bigr]^\top
      \in \mathbb{R}^6.
      \label{eq:legacy-state}
    \end{equation}
  \item \textbf{Inline the coupling logic} directly in the RHS function,
    computing the orifice flow rate and valve hysteresis in-place.
  \item \textbf{Hard-code} which tank receives the mission discharge, which
    tanks have vent edges, and what operating configuration (A/B/C) each tank
    is in.
  \item \textbf{Write separate thermal sub-routines} per tank, since there is
    no shared abstraction for the wall heat-transfer model.
\end{enumerate}

The resulting RHS function takes the schematic form shown in
\cref{lst:legacy-rhs}.

\begin{lstlisting}[language=Python,
    caption={Schematic legacy RHS for the two-tank CH2/CCH2 system.
             Valve state is maintained in a mutable closure variable.
             All indices, volumes, and parameters are hard-coded.},
    label=lst:legacy-rhs]
_valve_open = False          # mutable closure --- breaks parallelism

def ode_system(t, y, mission):
    m1, T1, Ts1, m2, T2, Ts2 = y        # hard-coded state layout
    rho1, rho2 = m1 / V1, m2 / V2

    # Thermodynamic properties (two separate blocks, one per tank)
    P1  = PropsSI('P',       'T', T1, 'D', rho1, 'hydrogen')
    P2  = PropsSI('P',       'T', T2, 'D', rho2, 'hydrogen')
    h1  = PropsSI('Hmass',   'T', T1, 'D', rho1, 'hydrogen')
    h2  = PropsSI('Hmass',   'T', T2, 'D', rho2, 'hydrogen')
    cv1 = PropsSI('CVMASS',  'T', T1, 'D', rho1, 'hydrogen')
    cv2 = PropsSI('CVMASS',  'T', T2, 'D', rho2, 'hydrogen')
    a1  = PropsSI('d(P)/d(T)|D', 'T', T1, 'D', rho1, 'hydrogen')
    a2  = PropsSI('d(P)/d(T)|D', 'T', T2, 'D', rho2, 'hydrogen')

    # Coupling: orifice from Tank 1 to Tank 2
    global _valve_open
    if _valve_open and P2 >= P_CLOSE:
        _valve_open = False
    elif not _valve_open and P2 <= P_OPEN:
        _valve_open = True

    mdot_12 = 0.0
    if _valve_open and P1 > P2:
        mdot_12 = Cd * A_ORIFICE * math.sqrt(2 * rho1 * (P1 - P2))
        mdot_12 = min(mdot_12, MAX_FLOW_12)

    # Mission discharge --- hard-coded to Tank 2
    mdot_disc = np.interp(t, mission['time'], mission['flow_rate'])

    # Mass balances
    dm1 = -mdot_12
    dm2 =  mdot_12 - mdot_disc

    # Energy balances
    # (outflow h_e == h_i, so enthalpy term vanishes for Tank 1 outflow)
    Q1, Q1_ext = wall_heat_transfer_1(T1, Ts1)
    Q2, Q2_ext = wall_heat_transfer_2(T2, Ts2)

    dT1 = ((T1/rho1)*a1*dm1 + Q1) / (m1*cv1)
    dT2 = (mdot_12*(h1 - h2) + (T2/rho2)*a2*(dm2) + Q2) / (m2*cv2)

    dTs1 = (Q1_ext - Q1) / (ms1 * cs1)
    dTs2 = (Q2_ext - Q2) / (ms2 * cs2)

    return [dm1, dT1, dTs1, dm2, dT2, dTs2]
\end{lstlisting}

\paragraph{Key structural limitations.}
\begin{itemize}
  \item \textbf{Non-extensible state vector.}  Adding a third tank (e.g.\ a
    LH2 storage for the return leg) requires expanding
    \cref{eq:legacy-state} to $\mathbb{R}^9$, rewriting every index reference
    in \texttt{ode\_system}, and manually introducing all new edge terms.

  \item \textbf{Inlined topology.}  The coupling graph ($1\to2$, discharge on
    $2$, vent on both) is embedded as flow-of-control in the RHS function.
    Swapping the mission to Tank~1, or adding a two-way coupling, requires
    code surgery rather than a parameter change.

  \item \textbf{Mutable valve state.}  The hysteresis flag \texttt{\_valve\_open}
    is a mutable closure variable.  This prevents the RHS from being called
    by a stiff solver at arbitrary time points (e.g.\ for Jacobian estimation)
    without corrupting the valve state.

  \item \textbf{Fixed-step integration.}  The legacy framework used a
    first-order Euler method with a fixed \SI{60}{\second} timestep.  For the
    CCH2 tank, whose transient at temperature \SI{53}{\kelvin} evolves on a
    much shorter timescale during rapid discharge events, this timestep
    introduces significant integration error.

  \item \textbf{No peripheral components.}  Conditioning the fluid stream
    (HEX, compressor, pressure reducer) between two nodes has no
    first-class representation; it would have to be inlined into the coupling
    term of the RHS.

  \item \textbf{Configuration B: discharge throttling.}  Under the legacy
    approach, reaching the minimum pressure floor caused the discharge rate
    to be algebraically throttled to whatever the ambient heat leak could
    sustain --- introducing a fictitious upper bound on mission flow.  This
    is physically incorrect when a heat exchanger or fuel-cell waste heat
    can compensate the pressure drop.
\end{itemize}

% -----------------------------------------------------------------------------
\subsection{New Paradigm: Graph-Assembled ODE}
\label{sec:comparison:new}
% -----------------------------------------------------------------------------

In the graph-based framework, the same physical system is declared as a network
in a YAML configuration file.  The ODE is never written by hand; it is
assembled automatically from the graph topology by the \texttt{TankSystem}
class.

\subsubsection{Network graph for the CH2/CCH2 case}

\Cref{fig:network-ch2-cch2} illustrates the directed graph
$G = (\Nodes, \Edges)$ for this topology.

\begin{figure}[ht]
\centering
\begin{tikzpicture}[
    tank/.style  = {draw, rounded corners=4pt, fill=blue!10, minimum width=2.8cm,
                    minimum height=1.0cm, align=center, font=\small\bfseries},
    sink/.style  = {draw, dashed, rounded corners=4pt, fill=gray!10,
                    minimum width=2.0cm, minimum height=0.8cm,
                    align=center, font=\small},
    arr/.style   = {-{Latex[length=6pt]}, thick},
    arr2/.style  = {-{Latex[length=6pt]}, thick, dashed, gray},
    lbl/.style   = {font=\footnotesize, midway},
    node distance=1.8cm,
  ]
  % Nodes
  \node[tank]  (T1)                        {Tank 1\\CH2\\800 bar / 288 K};
  \node[tank]  (T2) [right=4.5cm of T1]    {Tank 2\\CCH2\\400 bar / 53 K};
  \node[sink]  (FC) [below=2.0cm of T2]    {fuel cell\\(sink)};
  \node[sink]  (AT1)[above=1.6cm of T1]    {atmosphere\\(vent sink)};
  \node[sink]  (AT2)[above=1.6cm of T2]    {atmosphere\\(vent sink)};

  % Coupling edge
  \draw[arr] (T1.east) -- node[above, font=\footnotesize, align=center]{%
    \texttt{pressure\_compensation}\\$d=\SI{5}{\milli\metre}$,\;
    $P_\text{open}=\SI{16}{\bar}$} (T2.west);

  % Discharge edge
  \draw[arr] (T2.south) -- node[right, font=\footnotesize, align=center]{%
    \texttt{discharge}\\ATR\,72 profile} (FC.north);

  % Vent edges
  \draw[arr2] (T1.north) -- node[left,  font=\footnotesize]{\texttt{vent}} (AT1.south);
  \draw[arr2] (T2.north) -- node[right, font=\footnotesize]{\texttt{vent}} (AT2.south);
\end{tikzpicture}
\caption{Network graph $G$ for the coupled CH2/CCH2 system.  Solid arrows are
active edges; dashed arrows are vent edges that activate only in
Configuration~C.  All parameters live in the YAML file --- no physics code
is touched when the topology changes.}
\label{fig:network-ch2-cch2}
\end{figure}

\subsubsection{YAML network declaration}

The complete topology is declared in the \texttt{network} block of the scenario
YAML, shown in \cref{lst:yaml-cch2}.  Parameters for the valve model, the
mission profile, and each tank's initial/operating conditions are co-located in
a single self-describing file.

\begin{lstlisting}[language=Python,
    caption={Abbreviated YAML network block for the CH2/CCH2 case.},
    label=lst:yaml-cch2]
network:
  nodes:
    - node_id: 1
      type: tank
      fluid: CH2
      initial_conditions: {pressure: 80000000, temperature: 288.0}
      operating_limits:   {minimum_pressure: 1500000,
                           venting_pressure: 85000000}

    - node_id: 2
      type: tank
      fluid: CcH2
      initial_conditions: {pressure: 40000000, temperature: 53.25,
                           density: 80.0}
      operating_limits:   {minimum_pressure: 1500000,
                           venting_pressure: 45000000}

  edges:
    - edge_id: ch2_cch2_pressure_compensation
      connection_type: pressure_compensation
      from_node: 1
      to_node: 2
      activation_conditions:
        pressure_open_bar:  16.0
        pressure_close_bar: 30.0
        valve_time_constant_s: 0.5
      flow_physics:
        orifice_flow:     {orifice_diameter_m: 0.005}
        safety_limits:    {max_flow_rate_kg_s: 0.056}

mission:
  type: discharge
  profile: csv
  parameters:   {csv_file: .../atr72_fuel_flows.csv}
  assigned_to_node: 2
\end{lstlisting}

\subsubsection{Generic ODE assembly}

Given the graph $G$, the \texttt{TankSystem} assembler applies
\cref{eq:mass-balance,eq:energy-balance,eq:solid-temp} term-by-term for every
tank node by iterating over its incident edges.  For Tank~1
($\Edges_1^+ = \varnothing$, $\Edges_1^- = \{e_\text{coup}\}$):
\begin{align}
  \dot{m}_1 &= -\mdot_{e_\text{coup}}, \label{eq:new-mass1} \\[4pt]
  m_1 c_{v,1}\,\dot{T}_1 &=
    \underbrace{0}_{\substack{\text{no inflows,}\\\text{outflow}\;h_e=h_1}}
    + \frac{T_1}{\rho_1}\,\alpha_1\,(-\mdot_{e_\text{coup}})
    + \dot{Q}_{s,1}. \label{eq:new-energy1}
\end{align}
For Tank~2 ($\Edges_2^+ = \{e_\text{coup}\}$,
$\Edges_2^- = \{e_\text{disc},\, e_\text{vent}\}$):
\begin{align}
  \dot{m}_2 &= \mdot_{e_\text{coup}} - \mdot_\text{disc} - \mdot_\text{vent},
    \label{eq:new-mass2} \\[4pt]
  m_2 c_{v,2}\,\dot{T}_2 &=
    \mdot_{e_\text{coup}}\bigl(h_1 - h_2\bigr)
    + \frac{T_2}{\rho_2}\,\alpha_2\,\dot{m}_2^{\text{net}}
    + \dot{Q}_{s,2}
    + \dot{Q}_{HX,2}^{[\text{Config B}]}, \label{eq:new-energy2}
\end{align}
where $\dot{Q}_{HX,2}^{[\text{Config B}]}$ is non-zero only when Tank~2
operates in Configuration~B (\cref{sec:dae:configB}).  Equations
\cref{eq:new-mass1,eq:new-energy1,eq:new-mass2,eq:new-energy2} are exactly
equivalent to what the legacy developer would have written by hand in
\cref{lst:legacy-rhs} --- but here they are produced algorithmically with no
manual index tracking.

The coupling flow rate is computed by the \texttt{PressureCompensationValve}
subclass of \texttt{InterTankCoupling}, which encapsulates the orifice model
and hysteresis band:
\begin{equation}
  \mdot_{e_\text{coup}} =
  \begin{cases}
    \min\!\Bigl(C_d\,A_\text{orif}\sqrt{2\rho_1\max(0, P_1-P_2)},\;
               \mdot_{\max}\Bigr)
      & \text{valve open and } P_1 > P_2, \\
    0 & \text{otherwise},
  \end{cases}
  \label{eq:orifice-new}
\end{equation}
where the valve-state transition follows a hysteresis band on the \emph{target}
tank pressure ($P_\text{open} = \SI{16}{\bar}$, $P_\text{close} = \SI{30}{\bar}$),
and $A_\text{orif} = \pi d^2/4$ with $d = \SI{5}{\milli\metre}$.

The valve object carries no mutable state visible to the solver; the hysteresis
flag is updated at the beginning of each RHS call from the current state vector,
making the function pure enough for LSODA's internal Jacobian estimates.

% -----------------------------------------------------------------------------
\subsection{Advantages of the Graph-Based Paradigm}
\label{sec:comparison:advantages}
% -----------------------------------------------------------------------------

\paragraph{1 --- Topology-independent RHS assembler.}
The \texttt{TankSystem} assembler loops over nodes and their incident edges; the
$3N$-dimensional state vector and all ODE terms are constructed automatically.
Changing from a two-tank CH2/CCH2 system to a three-tank LH2/GH2/CCH2 system
requires editing only the YAML \texttt{network} block --- no RHS code is
modified.

\paragraph{2 --- Separation of model from scenario.}
All scenario parameters --- initial conditions, operating limits, mission
profiles, valve thresholds --- reside in a single YAML file.  The physics code
(\texttt{isochoric\_dynamic\_models.py}, \texttt{inter\_tank\_coupling.py})
is scenario-agnostic.  This separation enables:
\begin{itemize}
  \item version-controlled, human-readable scenario records;
  \item parametric sweeps by scripting YAML generation, without touching
    simulation code;
  \item reproducible baselines (the YAML file is a complete self-describing
    specification of the run).
\end{itemize}

\paragraph{3 --- First-class peripheral component chains.}
Attaching a heat exchanger and compressor to a coupling edge is a two-line YAML
addition (\texttt{components} list on the edge).  The assembler inserts the
component chain between the source node and the downstream energy balance.
Under the legacy paradigm such conditioning would require either (a) inlining
the component equations into the RHS function, tightly coupling the model to
that specific scenario, or (b) pre-computing conditioning results in a separate
script and interpolating them at run-time.

\paragraph{4 --- Physically correct Configuration~B.}
The legacy approach throttled the discharge rate to the algebraic limit
$\mdot_\text{disc}^{\max}$ --- the maximum sustainable by the ambient heat leak
alone.  This is physically incorrect whenever an on-board heat exchanger (e.g.\
fuel-cell waste heat) compensates the pressure drop.  The new approach instead
computes a compensating heat flux $\dot{Q}_{HX,i}$ (\cref{sec:dae:configB})
and keeps the discharge rate at the mission-prescribed value, maintaining the
system as a pure ODE throughout integration.

\paragraph{5 --- Adaptive stiff solver.}
The legacy framework used a fixed-step \SI{60}{\second} Euler method.  The
graph-based framework integrates with \textsc{lsoda} (via
\texttt{scipy.integrate.solve\_ivp}), which automatically selects between Adams
and BDF methods and adapts the step size.  For the CCH2 tank, whose pressure
evolves rapidly at \SI{53}{\kelvin} during coupling transients,
adaptive stepping reduces total RHS evaluations by an order of magnitude
compared with Euler at the same accuracy.

\paragraph{6 --- Valve modularity.}
Each valve archetype (\texttt{PressureCompensationValve},
\texttt{ProportionalSplitCoupling}, \texttt{PressureTriggeredDischarge}, \dots)
is an independent subclass of \texttt{InterTankCoupling}.  Swapping a
pressure-compensation valve for a proportional-split valve on the same edge
changes one YAML key (\texttt{connection\_type}) and requires no modifications
to the assembler or to other valve implementations.

\paragraph{7 --- Junction nodes enable complex topologies.}
The stateless \texttt{junction} node type (\cref{sec:network:junctions}) allows
flow-splitting and flow-merging without introducing dynamic states.  A
pressure-fed LH2 $\to$ GH2 buffer $\to$ mixer $\to$ fuel cell topology
(\cref{sec:network:yaml}) would require a 12-state monolithic RHS in the legacy
paradigm; in the new paradigm it is declared as five YAML edges.

\bigskip
\noindent\textbf{Summary.}  \Cref{tab:comparison} collects the key differences.

\begin{table}[ht]
\centering
\caption{Side-by-side comparison of the legacy and graph-based paradigms for
the CH2/CCH2 coupled system.}
\label{tab:comparison}
\begin{tabularx}{\linewidth}{@{} l X X @{}}
\toprule
\textbf{Aspect} & \textbf{Legacy} & \textbf{Graph-based} \\
\midrule
State vector
  & Manual index convention; must be expanded for each new tank
  & Assembled automatically: $3N$ states for $N$ tanks \\[4pt]
Topology declaration
  & Embedded as flow-of-control in \texttt{ode\_system}
  & YAML \texttt{network} block; no code change required \\[4pt]
Valve model
  & Inlined in RHS; mutable global state
  & \texttt{InterTankCoupling} subclass; pure function of current state \\[4pt]
Peripheral components
  & Must be inlined or pre-computed per scenario
  & Declarative \texttt{components} list on any edge \\[4pt]
Config~B
  & Discharge throttled to ambient heat-leak limit
  & Compensating heat flux; discharge at mission rate \\[4pt]
Integration
  & Fixed-step Euler (\SI{60}{\second})
  & LSODA (adaptive, stiff-aware) \\[4pt]
Extensibility
  & $O(\text{LOC})$ grows with each new tank or edge
  & $O(1)$: add a node/edge in YAML \\[4pt]
Reproducibility
  & Parameters scattered across Python constants
  & Fully described by a single YAML file \\
\bottomrule
\end{tabularx}
\end{table}
